\documentclass[11pt]{article}
\usepackage[utf8]{inputenc}
\usepackage{amsmath,amssymb}
\usepackage[margin=1in]{geometry}
\usepackage{hyperref}
\emergencystretch=3em

\title{Merging the log families of the reduced cutoff expansion}
\author{Claude, with Daniel Murfet}
\date{2026-09-07}

\begin{document}
\maketitle
\sloppy

\begin{abstract}
The reduced $d=2$ cutoff expansion listed three logarithmic families: the $u$-face logs, the
$v$-face logs, and the subtracted corner logs. All three are supported on the collision predicate
$(h_1+i+1)/k_1=(h_2+j+1)/k_2$, and for fixed $i$ (resp.\ $j$) at most one partner collides. Hence
the $v$-face family equals the corner family exactly and the reduced sum equals a merged sum with
ONE transferred log term $N^{-\alpha_i}(M[\alpha_i;0;c_{ij}/(k_1k_2)]\log N-M[\alpha_i;1;c_{ij}/(k_1k_2)])$
per collision $(i,j)$ (\texttt{reducedSum\_eq\_merged}). The proof is integrability-free. The analytic
cutoff theorems are restated in merged form. Zero \texttt{sorry}/\texttt{axiom}.
\end{abstract}

\section{The things formalised}
{\small
\begin{verbatim}
theorem transferTerm_ite (β α : ℝ) (n : ℕ) (P : Prop) [Decidable P] (g : ℝ → ℝ) (N : ℝ) :
    transferTerm β α n (fun s => if P then g s else 0) N
      = if P then transferTerm β α n g N else 0
theorem transferTerm_sum_ite_unique {ι : Type*} (β α : ℝ) (n : ℕ) (t : Finset ι)
    (Q : ι → Prop) [DecidablePred Q] (huniq : ∀ i ∈ t, ∀ j ∈ t, Q i → Q j → i = j)
    (g : ι → ℝ → ℝ) (N : ℝ) :
    transferTerm β α n (fun x => ∑ i ∈ t, if Q i then g i x else 0) N
      = ∑ i ∈ t, if Q i then transferTerm β α n (g i) N else 0
theorem uface_resonance_iff (h₁ h₂ k₁ k₂ i m : ℕ) (hk₂ : 0 < k₂) :
    ((m : ℝ) = (k₂ : ℝ) * ((((h₁ + i : ℕ) : ℝ) + 1) / k₁) - h₂ - 1)
      ↔ (((h₁ + i : ℕ) : ℝ) + 1) / k₁ = (((h₂ + m : ℕ) : ℝ) + 1) / k₂
noncomputable def mergedLog (β : ℝ) (h₁ h₂ k₁ k₂ M₁ M₂ : ℕ) (c : ℕ → ℕ → ℝ → ℝ)
    (N : ℝ) : ℝ :=
  ∑ i ∈ Finset.range M₁, ∑ j ∈ Finset.range M₂,
    if (((h₁ + i : ℕ) : ℝ) + 1) / k₁ = (((h₂ + j : ℕ) : ℝ) + 1) / k₂ then
      transferTerm β ((((h₁ + i : ℕ) : ℝ) + 1) / k₁) 1
        (fun s => c i j s / ((k₁ : ℝ) * k₂)) N
    else 0
noncomputable def mergedSum ... :=
  (u-face finite parts) + (v-face finite parts) + mergedLog
theorem reducedSum_eq_merged (β b : ℝ) (h₁ h₂ k₁ k₂ M₁ M₂ : ℕ) (hk₁ : 0 < k₁)
    (hk₂ : 0 < k₂) (a bj : ℕ → ℝ → ℝ → ℝ) (c : ℕ → ℕ → ℝ → ℝ) (N : ℝ) :
    reducedSum β b h₁ h₂ k₁ k₂ M₁ M₂ a bj c N = mergedSum β b h₁ h₂ k₁ k₂ M₁ M₂ a bj c N
theorem twoD_cutoff_analytic_merged ...   -- twoD_cutoff_analytic with mergedSum
theorem twoD_cutoff_amplitude_merged ...  -- twoD_cutoff_amplitude with mergedSum
\end{verbatim}
}

\section{Mathematics}

The three resonance predicates $m=k_2\alpha_i-h_2-1$ ($u$-face), $m=k_1\delta_j-h_1-1$ ($v$-face)
and $0=k_1\delta_j-(h_1+i)-1$ (corner) are all equivalent to the collision $\alpha_i=\delta_j$
with $\alpha_i=(h_1+i+1)/k_1$, $\delta_j=(h_2+j+1)/k_2$. The $u$-face log coefficient at $i$ is the
function $s\mapsto\sum_{j<M_2}[\alpha_i=\delta_j]c_{ij}(s)/(k_1k_2)$; since $j\mapsto\delta_j$ is
injective, at most one $j$ contributes, so the transferred term of this sum of indicators equals the
sum of indicator-weighted transferred terms (a two-case argument: a unique witness, or none, in which
case both sides vanish because \texttt{transferTerm} of the zero function is zero). The same applies
to the $v$-face family, which after \texttt{Finset.sum\_comm} coincides term by term with the corner
family (whose coefficient \texttt{faceLogCoeff} of \texttt{cornerData} is a single indicator). So
$U+V-C=U$, and $U$ is the merged family. On a collision $\alpha_i=\delta_j$, so the exponent may be
written either way. Astra~\#8's monomial tests hold for the merged form: with
$M[\alpha;\ell;f]=\int_0^\infty s^{\alpha-1}(\log s)^\ell e^{-\beta s^2}f(s)\,ds$ the collision term is
$N^{-\alpha}(M[\alpha;0;f]\log N-M[\alpha;1;f])/(k_1k_2)$ (sign fixed by
$\log(s/N)=\log s-\log N$), the two face finite parts add $(k_1+k_2)\log b\,M[\alpha;0;f]/(k_1k_2)$
at a collision with $b\ne1$, and off collisions the faces give
$b^{k_2(d-a)}N^{-a}M[a;0;f]/(k_1k_2(d-a))-b^{k_1(a-d)}N^{-d}M[d;0;f]/(k_1k_2(d-a))$.

\section{Paper-facing meaning}

This is the ``decisive bookkeeping lemma'' of Astra~\#8: the reduced sum is now in the normal form
$\sum_iN^{-\alpha_i}M[\alpha_i;0;U]+\sum_jN^{-\delta_j}M[\delta_j;0;V]+\sum_{\text{collisions}}N^{-\alpha_i}(M[\alpha_i;0;C]\log N-M[\alpha_i;1;C])$
from which the $d=2$ Taylor-tree statement (one polynomial $P_\mu$ of degree $\le1$ in $\log N$ per
pole $\mu$, coefficients independent of the cutoff) is a regrouping by exponent. Next: cutoff
independence of the finite parts and canonical $U_\alpha,V_\alpha,C_\alpha$, then the wrapper
\texttt{twoD\_taylor\_tree\_equal}.

\section{Lean notes}

\begin{itemize}
\item \texttt{transferTerm\_zero} already existed with a different meaning (log-degree $0$); the
zero-function lemma is \texttt{transferTerm\_zero\_fun}, following \texttt{logMoment\_zero\_fun}.
\item \texttt{if\_congr (hiff) (by ring) rfl} rewrites a decidable condition inside \texttt{ite}
together with its branch; \texttt{Finset.sum\_congr rfl fun j \_ => if\_congr ...} handles the
indicator sums without \texttt{simp} normalising the real equations.
\item Injectivity of $i\mapsto(h_1+i+1)/k_1$: \texttt{div\_left\_inj'}, \texttt{push\_cast},
then \texttt{exact\_mod\_cast} on a \texttt{linarith}-derived real equation.
\item The predicate equivalences are \texttt{eq\_div\_iff}/\texttt{div\_eq\_iff} followed by
\texttt{push\_cast} and \texttt{constructor <;> intro h <;> linarith}.
\item The whole unit compiled at the second attempt (the only failure was the name clash).
\end{itemize}

\end{document}
